\documentclass[11pt]{article}
\usepackage[utf8]{inputenc}
\renewcommand{\familydefault}{\sfdefault} % Makes Helvetica the default font
% set the page layout
\usepackage{geometry}
%set page spacing
\usepackage{setspace}
\usepackage{hyperref}

\usepackage{mathptmx}  % Times-like font
\geometry{
  a4paper,
  left=0.75in,
  right=0.75in,
  bottom=0.75in,
  top=0.75in
}
%empty the headers, footers, pagenumbers…
\pagestyle{empty}

% smaller itemize bullets
\renewcommand{\labelitemi}{\raisebox{0.07em}{\scriptsize$\bullet$}}

%Custom sectioning with secsty
\usepackage{sectsty}
\sectionfont{%                        
    \large % make sections smaller
    \fontfamily{ptm}\selectfont % change font family
    \sectionrule{0pt}{0pt}{-5pt}{1pt} % insert a thin rule
}

%% MACROS %%

% size of the boxes used to align text
\newlength{\spacebox}
\settowidth{\spacebox}{123456789}
% vertical space separator between entries
\newcommand{\sepspace}{\vspace*{1em}}

%name macro
\newcommand{\name}[1]{
    \Large%font size
    \fontfamily{ptm}\selectfont %font family
    %print name centered and bold
    \begin{center} \textbf{#1} \end{center}\par
    %back to normal size and font
    \normalsize
    \normalfont
}

\newcommand{\horizontalBox}[3]{
    \makebox[\linewidth][c]{
        \textsl{#1}\hfill$\bullet$\hfill
        \textsl{#2}\hfill$\bullet$\hfill
        \textsl{#3}
    }
        \noindent\rule{\linewidth}{0.4pt}% draws a thin line
}

\newcommand{\splitbox}[2]{
    \par\vspace{1em}
    \noindent\makebox[\linewidth][l]{
        \parbox[t]{0.25\linewidth}{\large \textbf{#1} \normalsize}
        \parbox[t]{0.75\linewidth}{#2}%
    }
}

%Subsection dated macro
\newcommand{\datedsubsection}[2]{%
  \subsection*{\normalsize \normalfont #1 \hfill #2}%
}

\newcommand{\skill}[2]{%
    \noindent
    \begin{minipage}[t]{0.45\textwidth} % Left column (skill name)
        \textsc{#1} % Small caps for skill title
    \end{minipage}%
    \hfill
    \begin{minipage}[t]{0.55\textwidth} % Right column (description)
        #2 % Multi-line description
    \end{minipage}%
    \vspace{0.8em}
    \par
}

\newcommand{\project}[4]{
    \noindent \textbf{#1}\normalsize \hfill #2
    \par
    \noindent #3
    #4
}

\newcommand{\job}[5]{
    \noindent \textbf{#1} \hfill #2
    \par
    #3
    \par
    #4
    \par
    #5
    \par
}

\begin{document}
    \begin{spacing}{1}
        \name{Connor Blaha}
        %Personal information
        \horizontalBox{blaha.s.connor@gmail.com}{330-606-0527}{Github: \url{https://github.com/SunCh1p}}
        %Summary
        \splitbox{Objective:}{ Computer Science graduate seeking a software development role or other tech related opportunities.}
        \splitbox{Education:}{BS Computer Science
            \par Graduated December 2025. \par Kent State University, Kent, Ohio \par GPA: 3.89}
        \splitbox{Skills:}{
            \textbf{Languages:} C++, Python, C\#, JavaScript, GDScript, HTML, CSS, SQL, Bash, LaTeX\\
            \textbf{Frameworks/Libraries:} React, Node.JS, Express, Flask, Pygame, Godot\\
            \textbf{Tools:} Git, VSCode, VisualStudio, OpenCode\\
            \textbf{Database Tools:} SQLite, MySQL}
        \splitbox{Coursework:}{Design Patterns, Algorithms, Human Interfaces and Computing, Introduction to Databases, Artificial intelligence, Introduction to Software Engineering, Introduction to Computer Graphics, Introduction to Game Design, Embedded Systems
        }
        \vspace{-2.5em}
        \section*{\Large \textbf{Projects}}
        \normalsize
        \project{Horror Game Prototype}{July 2026 - Present}{\textbf{Link:} \url{https://codeberg.org/SunCh1p/3D-Horror-Prototype}}
        {
            \begin{itemize}
                \setlength\itemsep{0.01em}
                \item Developing a first-person horror game in Godot 4.7 and GDScript, engineering modular player and enemy systems for extensible gameplay
                \item Implemented finite-state-machine enemy AI (idle/chase/attack) and player controller, ensuring predictable, debuggable behavior across gameplay states
                \item Engineered a signal-bus event system decoupling gameplay, UI, and audio subsystems to keep communication clean and maintainable; implemented procedural dungeon generation for replayable level variety
            \end{itemize}
            \par
        }
        \project{OpenGL Renderer}{March 2026 - July 2026}{\textbf{Link:} \url{https://codeberg.org/SunCh1p/OpenGL-GLFW}}
        {
            \begin{itemize}
                \setlength\itemsep{0.01em}
                \item Designed and built a 3D renderer in C++ with OpenGL and GLFW, structuring the pipeline for extensible multi-light rendering
                \item Implemented Phong shading with diffuse/specular maps and a GLSL material system to deliver realistic surface lighting
                \item Supported directional, point, and spot light casters, enabling multi-light scenes with consistent rendering performance
            \end{itemize}
            \par
        }
        \project{OpenGL Game}{August 2025 - December 2025}{\textbf{Link:} \url{https://github.com/SunCh1p/ComputerGraphicsClass}}
        {
            \begin{itemize}
                \setlength\itemsep{0.01em}
                \item Developed a 3D game in C++ and FREEGLUT, driving gameplay through a custom Entity Component System for clean separation of logic and rendering
                \item Implemented collision detection, animations, projectiles, model loading, and a scoring system, delivering a complete playable experience
            \end{itemize}
            \par
        }
        \project{Doc Reformed}{August 2025 - December 2025}{\textbf{Link:} \url{https://github.com/ch33s322/CS49999-Capstone}}
        {
            \begin{itemize}
                \setlength\itemsep{0.01em}
                \item Delivered a PDF-splitting desktop application in .NET and WPF as a capstone project, architected around a pipe-and-filter design for modular, extensible document processing
            \end{itemize}
            \par
        }
        \project{ItemQuest}{January 2025 - May 2025}{\textbf{Link:} \url{https://github.com/Item-Quest/TeamBravo}}
        {
            \begin{itemize}
                \setlength\itemsep{0.01em}
                \item Built a real-time multiplayer game with live scoring using Socket.IO, Flask, React, and TensorFlow.js, enabling simultaneous play with instant feedback
                \item Collaborated on a 10-person team using scrum methodology and a GitHub kanban board, ensuring on-time delivery within a single semester
            \end{itemize}
            \par
        }
    \end{spacing}
\end{document}
